% Additional Oct7 structural results and boundaries; not separate paper claims.
\section{Structural facts for effective-capacity cut payments}\label{noise:structure}
Throughout these appendices, $G$ is a finite directed acyclic $s$--$t$
multigraph with one fixed owner per edge, and every edge lies on an $s$--$t$
path. Let $F(c)$ be maxflow, $T=F(c)$, $D_e(c)$ the largest feasible flow
through edge $e$, $\mathcal C$ the structural inclusion-minimal cuts,
$K=|\mathcal C|$, and $L$ the longest path length. With independent uniforms
$U_e$ and a fixed report-independent cut priority, let $W$ minimize
$\sum_{e\in C}D_eU_e$ and put
\[
 z_e=D_e\Prob(e\in W),\qquad y_e(c)=\sup_{0\le t\le c_e}z_e(t,c_{-e}).
\]
Thus $y$ is the least own-report-monotone pointwise majorant of $z$.
For other noise laws the superscript specifies the law. Set $U(1)=1$ and,
for $K\ge2$, $U(K)=a_K+1/\log a_K$, where
$a_K\log(a_K/e)=\log K$ and $a_K>e$.
The usual uniform-noise tail calculation gives $z(E)\le U(K)T$:
normalized weights lie in $[0,1]$, some cut has total one, and
$\Prob(\sum a_eU_e\le1)\le(e/\rho)^\rho$ when $\sum a_e=\rho$.
These appendices distinguish bounds on a particular rule from bounds on
all eligible mechanisms. The latter sharp theorem is Theorem~\ref{prod:upper}.

\begin{lemma}[Sensitivity and a simultaneous flow]\label{noise:sensitivity}
Writing $c^t=(t,c_{-e})$ and $\delta_{eh}=D_h(c)-D_h(c^0)$, one has
\[
 D_h(c^t)=\min\{D_h(c),t+D_h(c^0)\},\quad
 0\le\delta_{eh}\le D_e(c).
\]
For every path antichain $A$, $\sum_{e\in A}\delta_{eh}\le D_h(c)$;
for all edges the analogous bound is $\sum_e\delta_{eh}\le LD_h(c)$.
Furthermore, if
$r_e=2\E[F(DU)-F((DU)_{-e},0)]$, then
$z\le r\le2z$, and a single feasible flow $\bar f$ at capacities $D$
satisfies $r\le\bar f$. Consequently $r(A)\le T$ and
$r(E)\le\min\{L,2U(K)\}T$.
\end{lemma}
\begin{proof}
On a DAG, the maximum possible flow through $h=(u,v)$ is the minimum of
$c_h$, the $s$--$u$ maxflow and the $v$--$t$ maxflow, omitting trivial
terminal terms. The prefix and suffix regions are disjoint by acyclicity,
so their attaining flows concatenate through $h$. Each relevant cut
contains the changing edge at most once. Its minimum therefore has form
$\min\{A+t,B\}$, which gives the displayed hinge on $[0,c_e]$;
for $h=e$ it is $\min\{t,D_e(c)\}$.
Choose one flow attaining $D_h$ and retain only its paths through $h$.
Deleting paths through $e$ proves $\delta_{eh}\le f_e$ simultaneously
for every $e$. Summing over an antichain or all edges gives the bounds.

At capacities $X=DU$, choose one maximum flow $f(X)$. Removing its paths
through $e$ proves $F(X)-F(X_{-e},0)\le f_e(X)$ simultaneously.
Conditionally on $U_{-e}$, the gain from setting the focal multiplier to
$u$ is a nonnegative increasing concave function $g(u)$ with $g(0)=0$.
Its endpoint is $g(1)=\min\{D_e,(B_e-A_e)_+\}$, where $A_e$ is the
cheapest containing-cut intercept and $B_e$ the cheapest excluding-cut
score. Conditional integration gives $\E g(1)=z_e$ and
$g(1)/2\le\int_0^1g(u)du\le g(1)$. Thus $z\le r\le2z$.
The flow $\bar f=2\E f(DU)$ conserves flow and obeys
$\bar f_e\le2D_e\E U_e=D_e$. A measurable selection is obtained from
the finitely many affine LP bases by lexicographic tie-breaking.
Path decomposition now gives the remaining claims.
\end{proof}
The hinge also gives a valid, though noncollective, gain estimate. Put
\[
 \Delta_e(t,U)=\max_{C\ni e}\sum_{h\in C\setminus\{e\}}
                         (\delta_{eh}-t)_+U_h.
\]
The excluding score cannot increase under an underreport, and the containing
intercept decreases by at most $\Delta_e$. The clamp is increasing and
one-Lipschitz, whence
\[
 0\le y_e-z_e\le\sup_{0\le t\le D_e}\E\min\{t,\Delta_e(t,U)\}.
\]
Reports above $D_e$ preserve the feasible-flow set. If no excluding cut
exists, the excess is zero. An expectation of a maximum cannot be replaced
by a maximum of expectations; no sharp aggregate gain bound follows here.

\begin{proposition}[Antichains and a conditional reduction]\label{noise:antichain}
For incomparable edges $e,f$, $y_e+y_f\le R(\{e,f\})$, where $R(A)$ is
the maximum flow on paths meeting the antichain $A$. For a complete parallel
bundle $X$ from $u$ to $v$, with total capacity $q$ and external prefix/suffix
bottleneck $B$, one has $y(X)\le\min\{q,B\}$.
If an inclusion-minimal antichain $A$ violates $y(A)\le R(A)$, then
$|A|\ge3$ and the simultaneous profile setting $c_e=y_e(c)$ for $e\in A$
preserves $R(A)$. At that profile every cut meeting $A$ costs strictly more
than $R(A)$, and there is a common unchanged minimum cut disjoint from $A$.
This reduction does not assert existence or impossibility of such an $A$.
\end{proposition}
\begin{proof}
The union $G_A$ of paths meeting $A$ has no path avoiding $A$: vertices before
and after $A$ cannot overlap, and an edge crossing between those regions
must belong to $A$, by acyclicity and the antichain property.
Fix $e\in A$, put $R=R(A)$ and $R_-=R(A\setminus\{e\})$. Under an
own underreport, if the effective focal capacity is at most $R-R_-$ its
payment is at most that quantity. Otherwise every containing cut in $G_A$
has capacity greater than $R$, since its other edges block an $R_-$ flow.
A minimum cut $C_0$ of $G_A$ excludes $e$ and has capacity at most $R$.
It blocks every path through $e$, so the endpoint flow marginal defining
$z_e$ is at most its random capacity. Taking its mean gives $R/2$. Hence
\[
 y_e\le\min\{D_e,\max(R-R_-,R/2)\}.                 \tag{A}
\]
For two edges, $R_-=D_f$. If either $D_f\le R/2$, then
$y_e\le R-D_f$ and $y_f\le D_f$; if both effective capacities exceed
$R/2$, each payment is at most $R/2$.

For a parallel bundle, $D_i=\min(c_i,B)$. External feasible flows depend
on its reports only through $\min(q,B)$. If $q\le B$, use $y_i\le c_i$.
Otherwise write $W=\sum_iD_i$. Every minimal cut contains the whole bundle
or none, so $z_i=D_ip$ for a common $p$; the simultaneous flow lemma gives
$Wp\le B$. With $s_i=q-c_i$ and $\delta_i=(B-s_i)_+$, underreports
$t\ge\delta_i$ keep every other effective cap fixed, and conditional
clamping implies payment at most $z_i(c)$. Smaller reports pay at most
$\delta_i$. Thus $y_i\le\max(z_i,\delta_i)$. If $\delta_i>0$, then
$W=D_i+s_i$ and $D_i\ge B-s_i$, giving
$\delta_iW\le BD_i$. Summing $y_i\le BD_i/W$ proves the bundle claim.

For a minimal violating $A$, write $Y(B)=\sum_{e\in B}y_e$ and
$\delta_e=R(A)-R(A\setminus\{e\})$. Minimality gives $y_e>\delta_e$.
The own-report maxflow hinge in $G_A$ is
$\min\{R(A),R(A\setminus\{e\})+t\}$, so any witness with
$z_e(c^t)>\delta_e$ preserves $R(A)$. No original minimum cut can contain
$e$, since such a cut would give $c_e=\delta_e$. Thus every original
minimum cut $C^*$ is disjoint from all of $A$.
For a cut $C$ let $B=C\cap A$. Its unchanged non-$A$ capacity is at least
$R(A\setminus B)$, by a flow meeting $A\setminus B$ and avoiding $B$.
At the joint $y$-profile its capacity, if $B\ne\varnothing$, is at least
$Y(B)+R(A\setminus B)\ge Y(A)>R(A)$. Cuts avoiding $A$ retain their
capacities and $C^*$ stays tight. It is saturated by a through-$A$ flow,
so its full-network effective capacities remain its original reports.
Near-maximizing separate witnesses suffice; no attainment is needed.
\end{proof}

\begin{proposition}[Cut height]\label{noise:height}
For $h(C)=\max_P|P\cap C|$ and $H=\sum_Ch(C)$,
\[
 h(C)\le\left\lfloor\frac{\sqrt{8K+17}-3}{2}\right\rfloor,
 \qquad H=O(K^{3/2}).
\]
The exponent for $H$ is attained in order by the classical Fibonacci DAG:
vertices $0,\ldots,2r-1$ and arcs of lengths one and two have
$K=r^2$ and $H=r(r+1)(2r+1)/6$.
\end{proposition}
\begin{proof}
For a minimal cut $C$, let $R$ be the vertices reachable after deleting it.
Minimality gives, for each $e\in C$, a private path meeting $C$ only in $e$;
its prefix lies in $R$ and suffix outside $R$, so $C=\delta^+(R)$.
Fix a path $P$ meeting $C$ in $e_i=(a_i,b_i)$, $i=1,\ldots,r$, and a
topological order $\tau$. For $i\le j$, the vertex set
\[
 S_{ij}=\{v:\tau(v)\le\tau(a_i)\}
       \cup\{v\in R:\tau(v)<\tau(b_j)\}
\]
has boundary trace on $P$ exactly $e_i,\ldots,e_j$. Each corresponding
private path meets that boundary only at its own edge, so every minimal
subcut contains the full interval. Distinct intervals give distinct cuts.
Every edge outside $C$ on $P$ also has a minimal cut with singleton $P$ trace,
obtained by a topological-prefix boundary. If $P$ has length $\ell$,
$K\ge\ell+r(r-1)/2$; successive exits from $R$ imply $\ell\ge2r-1$.
This proves the stated finite bound.

For the Fibonacci DAG, a canonical cut partition in vertex order has word
$S^a(TS)^bT^c$ with $a,c\ge1$: two consecutive $T$'s prevent a later
reachable $S$, while two consecutive $S$'s after the first $T$ prevent its
private suffix to the sink. Conversely such words make every $S$ reachable
within $S$ and every $T$ able to reach the sink within $T$, proving minimality.
The height is $h=b+1$, achieved by the backbone path. There are
$2r-2h+1$ such cuts, giving both sums. The graph and quadratic cut count
are established; no priority is asserted for them or the height comparison.
\end{proof}

\begin{proposition}[Serial blocks and failure of a maximum-module bound]
\label{noise:series}
If an edge is replaced by $k$ serial owners with no internal attachments,
their total envelope payment is at most their common effective capacity.
Nevertheless no constant multiple of the largest standalone module payment
bounds envelope payment after arbitrary series composition.
\end{proposition}
\begin{proof}
Let $q$ be the least report in the serial block. All its effective capacities
are a common $d(q)$, and external capacities depend only on $q$. Minimal cuts
use either no block edge or one arbitrary member. The smallest of its $k$
uniforms selects a uniformly distributed member, independently of that
minimum and all external uniforms. Thus every member has base payment
$Z(q)/k$, where $Z(q)=d(q)\Prob(W\text{ meets the block})$.
Each member's underreports realize exactly bottlenecks $q'\in[0,q]$;
therefore the sum of their envelopes is $\sup_{q'\le q}Z(q')\le d(q)$.
The collapsed noise is $\operatorname{Beta}(1,k)$, not uniform.

For the second claim, fix $m\ge2$, $r=4^m m!$, $M=16r$. Each of $r$
modules has two serial bundles: $M$ owners of capacity $1/M$ and $m$
owners of capacity one. Let their uniform scores be $L_i,H_i$.
Standalone payment is $1+(m-1)\Prob(H_i<L_i)\le1+(m-1)/m!\le2$.
In the composition, some $H_i\le1/4$ has probability at least $1/2$,
whereas $\Prob(\min_iL_i<1/4)\le r\,4/(3M)=1/12$.
A heavy cut therefore wins with probability at least $5/12$, giving payment
at least $5m/12$. The base equals its envelope at these profiles: underreporting
a light owner leaves all other owners of its containing cut unclipped, as
the resulting grand flow stays at least $1/M$; underreporting a heavy owner
leaves its pool total at least one. The containing intercept is fixed and
the other caps are nondecreasing, so conditional clamping proves monotonicity.
For general serial modules, $K=\sum_iK_i$ and global caps are
$D_{i,e}=\min\{d_{i,e},\min_jT_j\}$; this bookkeeping alone supplies no
maximum-module estimate.
\end{proof}

\section{What fixed noise and Bernoulli conditioning can achieve}\label{noise:laws}
\begin{theorem}[Conditioned-zero binary production]\label{noise:bernoulli}
Consider finitely many nonempty resource bundles $P$, each consuming one
unit of its resources and producing one unit of value. Each resource has
one owner. Let $V$ be the vertices of
$\{w\ge0:\sum_{e\in P}w_e\ge1\ \forall P\}$ and $K=|V|$.
Let $D_e$ be maximum feasible resource usage and $w^*$ the fixed-priority
minimizer of $\sum_eD_eB_ew_e$, where $\Prob(B_e=0)=q\in(0,1)$ independently.
Then
\[
 p_e=D_e\E[w_e^*\mid B_e=0]
\]
is own-report monotone, covers every coalition, and for every $\lambda>0$
obeys
\[
 p(E)\le\frac{T(\log K+\lambda+1)}
 {q[-\log(q+(1-q)e^{-\lambda})]}.
\]
For fixed $q=e^{-1}$ this is $(e+o(1))T\log K$. In the network subclass,
optimizing one common $q=q(K)$ cannot improve the leading constant $e$.
This class bound is superseded as an unrestricted existence bound by
Theorem~\ref{prod:upper}, but its formula and class limitation are distinct.
\end{theorem}
\begin{proof}
LP duality gives $F(c)=\min_{w\in V}c\cdot w$ and $0\le w_e\le1$:
a vertex coordinate exceeding one could be perturbed both ways without
violating an active constraint. Nonempty activities bound every primal
variable. Effective usage satisfies $D_e\le\min(c_e,T)$ and preserves
the entire feasible polytope, also after any binary mask. Convex mixtures
of feasible plans at own reports zero and $u$ imply, for $0<t\le u$,
\[
 D_h(t)\ge D_h(0)+(t/u)[D_h(u)-D_h(0)].
\]
Thus $D_h(t)>D_h(0)$ is independent of the positive own report. With $B_e=0$
the minimum masked value is independent of that report. A vertex remains
optimal precisely when it was optimal at zero and every coordinate with
$B_hw_h>0$ has no effective-usage increase. Hence the whole minimizing
family and its fixed-priority selection are invariant on positive own
reports. Multiplication by nondecreasing $D_e$ proves incentives, including
zero where the payment vanishes.
If $w_0,w_1$ are selected with $B_e=0,1$, their two optimality inequalities
give $w_{0e}\ge w_{1e}$ when $D_e>0$. Therefore for
$z_e^q=D_e\E w_e^*$, one has $z_e^q\le p_e\le z_e^q/q$.
Every dual vector covers every coalition under masked effective capacities,
so the same is true of $z^q$ and $p$.

For $a_e=D_ew_e/T\in[0,1]$ and $\rho_w=\sum a_e$, convexity gives
$\log(q+(1-q)e^{-\lambda a_e})\le-\alpha a_e$, where
$\alpha=-\log(q+(1-q)e^{-\lambda})$. Some unmasked objective equals $T$,
so the winning masked objective is at most $T$. Chernoff and a union over
vertices yield $\Prob(\rho_{w^*}>s)\le\min(1,Ke^{\lambda-\alpha s})$.
Integrating gives $\E\rho_{w^*}\le(\log K+\lambda+1)/\alpha$.
Take $q=e^{-1}$ and $\lambda=\log(2+\log K)$ for the asymptotic bound.

In graphs, $V$ consists exactly of minimal-cut vectors: maxflow/mincut
identifies the cut dominant with the path-cover polyhedron by separation
with nonnegative objectives; each minimal cut is uniquely exposed by
cost one on itself and a sufficiently large cost outside. With directed
cycles, define usage by path packings, excluding circulations. The same
binary theorem also covers unit-value splittable multicommodity paths,
where dual vertices can be fractional and $K$ is not a cut count.

For the sharp constant, the network identity
\[
 p^q(E)=q^{-1}\E\sum_{e\in W}D_e(1-B_e)                 \tag{B}
\]
is immediate by conditioning. With $K$ serial two-owner pools, each report
$1/2$, $T=1$ and
$p^q(E)=[1-\tfrac12(1-q^2)^K-\tfrac12(1-q)^{2K}]/q
\ge K/(1+2Kq)$, by $(1-q)^{-2K}\ge1+2Kq$.
Put $N=K-1$, $A=\log N\ge16$. If $q\ge A^{-2}$, write
$h=\log(1/q)$, $m=\lfloor(A-\sqrt A)/h\rfloor$,
$R=\lceil\sqrt A/h\rceil$. Use $N$ serial pools of $m$ unit owners
and a final pool of $R$ owners of report $1/R$.
When the final score is positive and some heavy pool is all closed, the
winning cut has closed weight $m$. By (B),
\[
 p^q(E)\ge\frac m q(1-q^R)[1-(1-q^m)^N]
 \ge e(A-\sqrt A-2\log A)(1-e^{-\sqrt A})(1-e^{-e^{\sqrt A}}),
\]
since $Nq^m\ge e^{\sqrt A}$, $q^R\le e^{-\sqrt A}$ and
$q\log(1/q)\le1/e$. If $q<A^{-2}$ the two-owner family gives
$p^q(E)\ge A^2/3\ge eA$. The lower is uniform in $q$, proving class
optimality on multigraphs, without claiming its simple-DAG conversion.
\end{proof}

\begin{theorem}[Fixed iid-law classification]\label{noise:classification}
Fix a law $X$ on $[0,1]$, independent of graph size, and replace uniforms
in $z,y$ by iid $X$. On DAGs the worst envelope factor has these properties:
identically zero noise has no owner-independent bound already at $K=2$;
constant positive noise has exact worst factor $K$ on multigraphs;
a nondegenerate law on $\{0,c\}$ has order $\Theta_q(1+\log K)$;
and a nonconstant conditional positive law has a lower bound
$\Omega_X(K^{\eta_X}/\log K)$ for some $\eta_X>0$, even on simple DAGs.
Every nonzero law has an $O_X(K)$ upper bound. These statements concern
pointwise-majorant repairs of that fixed base, not arbitrary mechanisms.
\end{theorem}
\begin{proof}
Let $\mu=\E X>0$. For a cut $C$, set $w_h=D_h/T\le1$,
$d=\sum_Cw_h$ and $S=\sum_Cw_hX_h$. At a report $t=xD_e$, $0<x\le1$,
homothetic scaling of feasible flows gives $D_h(c^t)\ge xD_h(c)$.
Selection of $C$ implies $xS\le1$. For $u=1/x$, if $u\ge\mu d/2$,
$(d/u)\Prob(S\le u)\le2/\mu$; otherwise Hoeffding gives a bound
$d e^{-\mu^2d/2}\le2/(e\mu^2)$. Thus
\[
 t\Prob(W(c^t)=C)\le A_\mu T\frac{D_e}{D(C)},\qquad
 A_\mu=\max\{2/\mu,2/(e\mu^2)\}.
\]
Summing over containing cuts before taking suprema proves $y(E)\le A_\mu KT$.
Constant positive noise makes the same argument deterministic with constant
one. Equality is approached with $K$ serial two-owner pools of reports
$1/2$: lowering either owner makes its pool uniquely cheapest. On simple
DAGs $N$ serial diamonds give $K=4N$ and factor at least $N$, by lowering
an owner of each module's first-priority local cut.
Zero noise always selects the fixed first cut: two serial pools of $m$
owners, with reports one in that cut and $1/m$ in the other, give factor
$m$ at $K=2$. Subdivision into two-edge private paths gives a simple-DAG
lower of order $\log K$.
For $X\in\{0,c\}$, Theorem~\ref{noise:bernoulli} majorizes its envelope
and proves the logarithmic upper. For the lower, take
$N=\lceil q^{-m}\rceil$ pools of $m$ unit owners and a final unit singleton.
When the singleton is open and some heavy pool is closed, the winning
uncensored cut weight is $m$, so $z(E)\ge(1-q)(1-e^{-1})m$.
Subdivision of heavy edges replaces $q$ by $q'=1-(1-q)^2$ and gives
$K=N2^m+1$, still $m=\Theta_q(\log K)$ on simple DAGs.

It remains to handle a nonconstant conditional positive law. Let
$s=\Prob(X>0)$, $\nu=\E\min(X,X')$ and $d=\mu-\nu/s>0$;
strictness follows by conditioning on both draws being positive. Put
$a=(1/s+\mu/\nu)/2$, $b=d/(4\mu)$, $\delta=d/40$.
Then
\[
 \E\min(aX,(a+b)X')\le a\nu+b\mu=\mu-10\delta,
 \qquad \E\min(aX,mX')\uparrow a\mu s>\mu+20\delta.
\]
Choose $m_0$ making the second expectation at least $\mu+10\delta$;
integers $C\ge(\mu+10\delta)/(b\mu)$, $M\ge\delta^{-2}$; and
\[
 \gamma=\min\{32\delta^2/a^2,32\delta^2/(Cb^2)\},\qquad
 B\ge\max\{1,m_0,a+b,b/\delta,2/\gamma\}.
\]
For integer $r\ge4$ use $N=2^r$ modules in series, $m=Br$,
$k=Cm$, and a final $M$-owner pool of total one. In each module put
$e:s\to A$ of report one, $y_i:s\to B_i$ of report $a/m$,
$g_i:A\to B_i,h_i:B_i\to t$ of report two, and $k$ parallel
$Z:A\to t$ owners of report two. Under a separate focal report $q$,
$T=1$ and
$D_e=D_g=D_Z=q$, $D_y=a/m$, $D_h=\min(1,q+a/m)$.
These values follow from explicit paths; each raw module has flow $a+q>1$.
There are $2^{m+1}$ cuts per module, according as $e$ is included or
excluded, so $K=N2^{m+1}+1$.

Lower one focal report to $t=b/m$. Its cheapest containing score is
$tX_e+A_{\rm low}$, where
$A_q=\sum_i\min((a/m)X_{y_i},D_{h_i}X_{h_i})$.
The high and low means are at least $\mu+10\delta$ and at most
$\mu-10\delta$. Every excluding score is at least $q\sum_ZX_Z$,
whose low mean is at least $\mu+10\delta$.
The squared-range sums are $a^2/m$ and $Cb^2/m$. Hoeffding bounds each
wrong-side deviation to $\mu\pm6\delta$ by $e^{-\gamma m}$.
The final score is an average of $M$ draws and lies within $2\delta$
of $\mu$ except with probability $1/16$. A union over at most $2N$
other failures costs at most $1/8$. Since $t\le\delta$, every winning
cut contains the lowered focal owner on this event. Hence at the single
original profile $y(E)\ge Nb/(2m)$, proving a positive polynomial exponent.

For a simple DAG, choose the constants afresh for $V=\min(X_1,X_2)$,
whose conditional positive law is still nonconstant. Subdivide each
$y,h,Z$ and final edge into a private two-edge path, retaining $e,g$.
The primitive owner noise is still $X$, while those path scores use $V$.
The proof used only $X_e\le1$ and nonnegativity of $X_g$ for the remaining
edges. Now
$K=2^{(1+2B)r}+2^{[1+(C+2)B]r}+2^M$ and the same payment lower holds,
giving exponent $1/[1+(C+2)B]>0$. Strict cut-class separation makes all
arguments valid with atoms and any fixed tie priority.
\end{proof}

The classification fixes the law as $K$ grows. One particular variable-law
escape also fails: if $q=K^{-1/2}$ and
$X\sim(1-q)\delta_1+q\operatorname{Unif}(0,\epsilon)$ with
$0<\epsilon\le1/8$, use $K$ serial pools of reports $(3/4,1/4)$.
Lower each small owner separately to $1/8$. Its pool wins strictly when
the companion has low noise, the focal noise is one, and no other pool has
two low noises. Its score is at most $7/32$, while competitors cost at least
$1/4$. Effective capacities remain the reports. Summing gives
\[
 y(E)\ge(K/8)q(1-q)(1-q^2)^{K-1}\ge\sqrt K/64\quad(K\ge4).
\]
This does not classify arbitrary $K$-dependent noise laws.

\begin{proposition}[Uniform mixture of Bernoulli parameters]\label{noise:uniform-mixture}
The truthful covering rule $\bar p=\int_0^1p^q\,dq$ has worst-case
budget order $\Theta(K)$ on multigraphs.
\end{proposition}
\begin{proof}
Under a uniform random arrival order, before owner $e$ arrives select the
first cut minimizing the revealed $D$-weight and charge $D_e$ if it contains
$e$. Conditioning on its arrival time gives exactly $\bar p$.
At the last positive charge assigned to a particular cut, its revealed score
was at most $T$ before arrival and increases by at most $D_e\le T$.
All charges previously assigned to that cut are included in its revealed
score, even if the cut ceased to win and returned. Thus each cut is charged
at most $2T$, proving $\bar p(E)\le2KT$.
For the lower, take $K$ serial pools of $B^{K-1}$ owners, with $B\ge3$.
Enumerate pools in their fixed priority order. In pool $i$ give
$m_i=B^{K-i}$ owners report one and all others zero. Then $T=1$ and $D=c$.
Let $\tau_i$ be its first positive-owner arrival. If
$\tau_i>\max_{j<i}\tau_j$, cut $i$ is selected just before that arrival
and charges one. Since
$\Prob(\tau_i<\tau_j)=m_i/(m_i+m_j)\le B^{-(i-j)}$ for $j<i$,
\[
 \bar p(E)\ge K-\sum_i\sum_{d=1}^{i-1}B^{-d}
             \ge K(1-1/(B-1)).
\]
Zero owners remain part of the fixed topology. This construction has
$n=KB^{K-1}$ owners and does not refute an $O(\log n)$ bound.
\end{proof}

\section{Truthful repairs that lose their budget or coalition guarantee}
\label{noise:repairs}
\begin{proposition}[Deletion repairs can have linear cost]\label{noise:deletion}
Both of the following rules are truthful pointwise majorants of $z$, and
hence cover coalitions and majorize $y$:
\[
 p_e^{\rm pot}=2[\Psi(c)-\Psi(c^0)],\quad
 \Psi(c)=\E F(D(c)U),
 \qquad
 p_e^{\rm one}=\E\min\{D_e(c),(B_e(c)-A_e(c^0))_+\}.
\]
Here $A_e$ omits the focal term from containing cuts, and $B_e$ minimizes
over excluding cuts. If no excluding cut exists, set $p_e^{\rm one}=D_e$.
Each rule nevertheless has examples with total payment $\Omega(KT)$.
\end{proposition}
\begin{proof}
$\Psi$ is coordinatewise increasing, so its own deletion difference is
own-monotone. At fixed other uniforms, concavity of the focal flow gain
implies its average is at least half its endpoint. Since recomputing $D$
after deletion only decreases other capacities, $p^{\rm pot}\ge z$.
For the one-sided rule, $A_e(c^0)$ is independent of $c_e$, whereas $D_e$
and $B_e(c)$ are nondecreasing. Also $A_e(c^0)\le A_e(c)$, proving domination.

For the potential example, put a branch of $k$ unit serial owners in parallel
with a single owner of capacity $\beta=1/2$, followed by a final unit edge.
Then $T=1,K=k+1$. With $A_k=\min_{i\le k}U_i$,
$\Psi=\E\min(A_k+\beta U,V)$. Deleting any serial owner leaves effective
capacities $\beta$ on the other branch and final edge, so its deletion
potential is $\beta/3$. Meanwhile
$\Psi\ge\E\min(\beta U,V)=\beta/2-\beta^2/6$.
Each serial owner receives at least $\beta(1-\beta)/3=1/12$.
Indeed $A_k\to0$ and deletion of the parallel owner leaves potential
$1/(k+2)$, yielding
$p^{\rm pot}(E)=2[(k+2)\Psi-k/6-1/(k+2)]\sim k/12$.
The original envelope on this TTSP graph equals $z$ and has total at most
$3/2$, the largest effective weight of a selected cut.

For the one-sided example, put $k$ unit serial owners $X$ on $s\to A$,
add $b:s\to B$ of capacity $\beta\in(0,1)$, and unit edges
$a:A\to t,g:A\to B,h:B\to t$. Then $T=1+\beta$, $K=2k+2$; cuts are
$\{x_i,b\},\{x_i,h\},\{a,h\},\{a,b,g\}$.
Deleting $x_i$ leaves only $b+h$, so
$A_i(c^0)=\beta\min(U_b,U_h)$, while
$B_i(c)\ge m=\min(\beta U_b,U_h)$. Therefore
\[
 p_i^{\rm one}\ge\E[m-\beta\min(U_b,U_h)]
 =\beta/2-\beta^2/6-\beta/3=\beta(1-\beta)/6.
\]
For $k\ge2$, $B_i\le m+\min_{j\ne i}U_j$ and the expected error is
$1/k$. The four other payments sum to at most $3+\beta$, so
$p^{\rm one}(E)=k\beta(1-\beta)/6+O(1)$.
At $\beta=1/2$ the normalized asymptotic is $K/72$.
By contrast, when $x_i$ underreports, all serial effective caps become equal;
its selection requires the smallest of their uniforms. Hence $y_i\le1/k$
and $y(E)\le4+\beta$ on this family.
\end{proof}

\begin{proposition}[Common scales]\label{noise:scales}
Common clipping $a_e=\min(c_e,T)$ preserves coalition coverage and the
$U(K)T$ base budget but its uniform-cut payment need not be own-monotone.
The truncated scale envelope
\[
 P_e=\sup_{0\le r\le T}\min(c_e,r)
       \Prob\bigl(e\in W(\min(c,r)U)\bigr)
\]
is truthful and covering, but can cost at least $5KT/8$ with $O(K^7)$ owners.
\end{proposition}
\begin{proof}
Every feasible DAG flow has each coordinate at most its total, at most $T$,
so common clipping preserves all feasible flows. A minimum cut retains
weight $T$, proving the usual tail-budget premise.
For nonmonotonicity, use $x:s\to A$ of capacity one, $z:A\to t$ of
capacity three, $g:A\to B$ of capacity $\delta=1/10$, focal $h:B\to t$,
and $N$ parallel $s\to B$ owners of report $1/N$.
The four cuts are $\{x\}\cup Y$, $\{x,h\}$, $\{z,h\}$,
$\{z,g\}\cup Y$. For $h\in[3/4,1]$, $T=1+h$.
The focal payment is
\[
 w_h(h)=\E\min(h,Y+W_T),\qquad
 W_T=\min(\delta U_g,(U_x-TU_z)_+),\quad
 Y=N^{-1}\sum U_i.
\]
Direct tail integration gives
$\Prob(W_T>r)=(1-r/\delta)(1-r)^2/(2T)$ on $[0,\delta]$ and
$\E W_T=187/(8000T)$. At $h=3/4$, clamping loses at most
$\Prob(Y>13/20)\le1/2700$ when $N=10000$, by variance $1/(12N)$.
Ignoring the high-profile clamp only raises its payment. Thus
\[
 w_h(3/4)-w_h(1)\ge187/112000-1/2700=3929/3024000>0.
\]
A smaller exact certificate uses $N=6$. With
$F_N(u)=N^N\sum_{j=0}^N(-1)^j\binom Nj(u-j/N)_+^N/N!$ and
$L_N(u)=N^N\sum_{j=0}^N(-1)^j\binom Nj(u-j/N)_+^{N+1}/(N+1)!$,
the exact expression is
$w_h=1/2-L_N(1-h)+\int_0^\delta
(1-r/\delta)(1-r)^2F_N(h-r)/(2T)\,dr$.
Piecewise polynomial integration gives the positive difference
$587847649425223/914457600000000000$.

For a fixed external $r$, conditional clamping proves own-monotonicity;
the interval $[0,T]$ expands with the own report. The supremum is therefore
truthful and includes the covering rule at $r=T$.
For its lower bound, use $K$ serial bundles. Put $M=64K^4$, $N=256K^6$,
$b_j=3/4+j^2/(4K^2)$. Bundle $j$ has $N$ owners of report $b_j/N$ and
$n_j=M(K-j)$ owners of report one. Then $T=1$.
At scale $r_i=i/(2K^2M)$ all light owners are unclipped, and the clipped
bundle totals $g_j=b_j+n_jr_i$ satisfy
\[
 g_j-g_i=(j-i)^2/(4K^2),\qquad g_i=1-(K-i)^2/(4K^2)\ge3/4.
\]
Each random score has variance at most $1/(12N)+1/(48KM)$.
Chebyshev and a union bound with $\eta=1/(16K^2)$ give failure probability
at most $1/(12K)+1/12\le1/6$ for all deviations being less than $\eta$.
The mean gaps are at least $2\eta$, so bundle $i$ wins with probability
at least $5/6$. Summing its scale-$r_i$ payments over disjoint bundles
proves $P(E)\ge5K/8$.
The same obstruction can be made balanced: replace heavy reports in
$j<K$ by $(K+j)/(4K^2M)$. Every raw pool then totals one; for $i<K$
the clipped totals become $\min(1,g_j)$, preserve the stated gap, and have
no larger variance. This gives $P(E)\ge5(K-1)/8$.
\end{proof}

\begin{proposition}[Personal ceilings lose coverage]\label{noise:ceilings}
Let $B_e=\min_{C\not\ni e}c(C)$, allowing $\infty$, and use for owner
$e$ the clipped vector $a_h^{(e)}=\min(c_h,B_e)$ and payment
$p_e=a_e^{(e)}\Prob(e\in W(a^{(e)}U))$.
This is truthful and spends at most $(U(K)+1)T$, but no universal constant
multiple covers every coalition, even on TTSP graphs.
\end{proposition}
\begin{proof}
$B_e$ is independent of $c_e$, proving incentives by conditional clamping.
If $B_e=T$, its coordinate equals the common-clipped payment. If $B_e>T$,
every minimum cut contains $e$, so all such exceptional edges belong to one
minimum cut and their payments sum to at most $T$. The rest sum to at most
$U(K)T$. At $T=0$ the corresponding endpoint marginals vanish.

For a counterfamily, fix $q\ge4$, $m=q^2$, $\delta=q^{-2}$,
$N=m!q^m$. The initial component has parallel branches: an edge $e$ of
capacity $\delta$, and a serial path with first capacity $1-\delta$ and
$m$ further owners of report $m$. Follow it by $N$ serial bundles of $m$
unit owners each. Cuts are $\{e,a\}$, $\{e,h_i\}$ and those bundles;
$T=1,K=N+m+1$. Personal ceilings are $B_e=B_a=m$ and all others one.
Take $S$ to contain $e$ and one owner from every bundle, so $v(S)=\delta$.
With $Y=\min_{i\le m}U_{h_i}$,
$Z=\min_{j\le N}\sum_{k\le m}U_{jk}$ and
$X_r=\delta U_e+\min((1-\delta)U_a,rY)$, its payment is exactly
\[
 p(S)=\delta\Prob(X_m<Z)+\Prob(Z<X_1).
\]
The simplex CDF $z^m/m!$ for $z\le1$ gives
$\Prob(Z<1/(2q))\le2^{-m}$ and
$\Prob(Z>2/q)\le(1+2^m)^{-1}$.
On $X_m<Z\le2/q$, either $(1-\delta)U_a<2/q$ or $mY<2/q$.
Also $X_1\le\delta+Y$. Hence
\[
 \frac{p(S)}{v(S)}\le
 \frac{2q}{q^2-1}+\frac2q+\frac1{1+2^{q^2}}
 +q^2\left[2^{-q^2}+
  \left(1-\frac{q-2}{2q^2}\right)^{q^2}\right]
 =\frac{4+o(1)}q\longrightarrow0.
\]
Indeed the last power is at most $e^{1-q/2}$.
Stirling gives $\log K/\log\log K\sim3q^2/2$, so the necessary scalar
coverage repair is at least order $\sqrt{\log K/\log\log K}$.
This is a limitation of owner-dependent normalization, not a lower bound
on optimal mechanisms.
\end{proof}

\subsection{Failure of coordinatewise and cross-owner comparisons}
\begin{proposition}\label{noise:pointwise}
Neither $y_e\le Cz_e$ nor $y_e\le Cr_e$ holds for a universal constant
$C$. This already fails for DAGs of longest path four and cut height at
most two. The frozen marginal $r$ itself need not be own-monotone.
\end{proposition}
\begin{proof}
Use vertices $s,A,B_1,\ldots,B_m,q,t$: focal $x:s\to A$;
$y_i:s\to B_i$ of report $\epsilon=1/10$;
$g_i:A\to B_i,h_i:B_i\to q$ of report two;
$k$ parallel $Z:A\to q$ owners of report two;
$N$ direct $P:s\to q$ owners of report $(1-2b)/N$;
and $N$ final $R:q\to t$ owners of report $1/N$.
Whenever $m\epsilon\ge2b$ and $x\in[0,9/10]$, $T=1$ and
$D_x=x,D_y=\epsilon,D_g=D_Z=x,D_h=x+\epsilon$, with $D_P,D_R$ unchanged.
Cut families are $R$; $P\cup\{x\}$ plus one of $y_i,h_i$ per branch;
and $P\cup Z$ plus one of $h_i,\{y_i,g_i\}$ per branch.
Thus $K=2^{m+1}+1$, and private paths verify every listed cut is minimal.

Let $B=R-P$ now denote random score sums, so $\E B=b$ and
$\operatorname{Var}B\le1/(6N)$. Put
$A_i=\min(\epsilon U_i,(x+\epsilon)V_i)$ and
$J_x(s)=\E\min(x,(s-\sum_iA_i)_+)$.
If the random $Z$ score is at least one, the only competing excluding
threshold is $R$. Conditional integration and Lipschitzness therefore give
\[
 |z_x-J_x(b)|\le1/\sqrt{6N}+10^k/k!.
\]
For $b\le\min(\epsilon,x)$ the density of each $A_i$ on the relevant
simplex is $\alpha_x-\beta_xa$, with
$\alpha_x=10+1/(x+\epsilon)$, $\beta_x=20/(x+\epsilon)$. Thus
\[
 J_x(b)=\sum_{j=0}^m\binom mj\alpha_x^{m-j}(-\beta_x)^j
                      \frac{b^{m+j+1}}{(m+j+1)!}.       \tag{C}
\]
Set $b=1/(100m)$, $I=b^{m+1}/(m+1)!$, $Q=11^mI$,
$N=\lceil16/Q^2\rceil$, $k=2\lceil8/Q\rceil$.
At $x_0=1/10$, density is at least $15-1/m$, so
$J_{x_0}\ge(14/15)15^mI$; at $x_1=9/10$ it gives
$(54/55)Q\le J_{x_1}\le Q$.
The error is at most $Q/4$: the variance term is at most $Q/8$, and
$10^k/k!\le(200/k)^{k/2}\le2^{-k}\le1/k\le Q/16$ since $k\ge800$.
Consequently $z_{x_1}>0$ and
\[
 \frac{y_x(c^{x_1})}{z_x(c^{x_1})}
 \ge\frac{56}{75}(15/11)^m-\frac15\longrightarrow\infty.
\]
The comparison with $r$ fails because $r\le2z$.
An exact smaller factor-two certificate uses $m=3,b=1/20,N=100000,k=50$;
(C) gives $J_{x_0}-2J_{x_1}=7627/134400000$.
The convolution density bound implies $|J_x''|\le\alpha_x$, so centering
at $b$ improves the error to $\alpha_x/(12N)$ and yields
$z_{x_0}-2z_{x_1}\ge1161/44800000-10^{50}/50!>0$.

For nonmonotonicity of $r$, use the same graph with $m=1,b=1/20$,
$N=100000,k=20$, and compare $x_0=1/2,x_1=3/5$.
Conditional frozen payment is
$\ell_x(t)=0$ for $t\le0$, $2t-t^2/x$ for $0<t<x$, and $x$ thereafter.
With $A=\min(\epsilon U,(x+\epsilon)V)$, direct integration gives
\[
 \E\ell_x(b-A)=(\epsilon^{-1}+(x+\epsilon)^{-1})b^2
 -\frac{2b^3}{3\epsilon(x+\epsilon)}
 -\frac{(\epsilon^{-1}+(x+\epsilon)^{-1})b^3}{3x}
 +\frac{b^4}{6\epsilon(x+\epsilon)x}.
\]
Its low-minus-high value is $23/100800$. The averaged function has
second derivative $2f_A(s)-(2/x)\Prob(s-x<A<s)$, bounded in absolute value
by $24$, so centering error per profile is at most $2/N$.
Truncation error is at most $2^{20}/20!$ per profile. Therefore
\[
 r_x(x_0)-r_x(x_1)\ge23/100800-4/100000-2^{21}/20!>0.
\]
These are finite fixed-owner examples, not limit-only counterexamples.
\end{proof}

Even a constant through-antichain value and a common disjoint minimum cut
do not imply cross-substitutes. Here is an exact separating construction.
Take two source owners $e,h$ of reports $1,h$ in parallel, followed by two
owners of report two in parallel. In parallel with that whole branch put a
100-edge path with every report ten. Follow both by $N=10000$ parallel
owners of total $21/2$. The source antichain is $\{e,h,p_1\}$.
On the joint box $e\in[.99,1]$, $h\in[.1,1]$, $p_1\in[9.9,10]$,
every upstream cut costs at least $10.99$, so the final cut fixes both
$T$ and through-antichain rank at $10.5$.
At $e=1,p_1=10$, the isolated first branch gives
\[
 q_h=\Prob(U_e+hU_h<(1+h)(U_{b_1}+U_{b_2}))
     =5/6+h/[12(1+h)^2],\quad q_1-q_{.1}=27/1936.
\]
In the full graph the additional condition is that this winning score plus
$10\min_{j\le100}U_{p_j}$ is below the final score. Failure probability
is at most $(9/10)^{100}+49/(27N)<1/4000$, by Chebyshev below three for
the final score. Thus $z_e(h=1)-z_e(h=.1)>27/1936-1/2000>0$.
The graph is TTSP and its envelope is safe; this only rejects the proposed
comparison and does not contradict Proposition~\ref{noise:antichain}.

